\begin{solucion}
Dado que el grado de $x^p-2$ es $p$, el polinomio $f$ tiene $p$ raíces en su cuerpo de descomposición. Asi sean $\alpha_1, \alpha_2, \ldots, \alpha_p \in L$ las raíces de $f$, entonces, por definición de cuerpo de descomposición:
\begin{align*}
L &= K(\alpha_1, \alpha_2, \ldots, \alpha_p) \\
x^p - 2 &= c(\alpha_1 - x)(\alpha_2 - x) \cdots (\alpha_p - x) \in L[x],
\end{align*}
Aplicando el criterio de Eisenstein con $p=2$:
\[
a_{p}=1\not\equiv0\pmod2,\qquad
a_{p-1}=\dots=a_{1}=0\equiv0\pmod2,\qquad
a_{0}=-2\not\equiv0\pmod{2^{2}}.
\]
Concluimos que $f$ es irreducible sobre $\Q$. Ahora sea $\alpha \in L$ es una raíz de $f$, entonces:
\begin{align*}
\alpha^p - 2 =0 \quad &\implies \quad \alpha^p = 2 = 2 e^{2\pi i}\\
&\implies \quad \alpha_k = \sqrt[p]{2} e^{2\pi i k/p} = \sqrt[p]{2} \omega_k, \quad k=0,1,\ldots,p-1,\\
\end{align*}
Como $\omega_0 = 1$, entonces $\alpha_0 = \sqrt[p]{2}$ es una raíz real de $f$ luego $\left(x-\sqrt[p]{2}\right)\mid f(x)$ es decir:
\begin{align*}
f(x) &= \left(x-\sqrt[p]{2}\right)\left(x^{p-1} + \sqrt[p]{2} x^{p-2} + \left(\sqrt[p]{2}\right)^2 x^{p-3} + \ldots + \left(\sqrt[p]{2}\right)^{p-1}\right)\in \mathbb Q(\sqrt[p]{2})[x].
\end{align*}
Pues tiene coeficiente racionales que están en $\mathbb Q(\sqrt[p]{2})$ y claramente $\sqrt[p]{2} \in \mathbb Q(\sqrt[p]{2})$. Notemos que:
\begin{align*}
  Q(\sqrt[p]{2})= \mathbb Q[\sqrt[p]{2}] &= \{a_0 + a_1 \sqrt[p]{2} + a_2 \left(\sqrt[p]{2}\right)^2 + \ldots + a_{p-1} \left(\sqrt[p]{2}\right)^{p-1} : a_i \in \mathbb Q\}\\
\end{align*}
Por lo tanto:
\begin{align*}
  [\mathbb Q(\sqrt[p]{2}):\Q] &= p \quad \text{( definición de grado de una extensión)}\\
\end{align*}
Por otro lado dado que:
\begin{align*}
  \omega_k &= e^{2\pi i k/p} \\
  &= \left(e^{2\pi i/p}\right)^k \\
  &= (\omega_1)^k,
\end{align*}
Entonces $\mathbb Q(\omega_1, \ldots, \omega_{p-1}) = \mathbb Q(\omega_1)$, pues claramente $\mathbb Q(\omega_1) \subseteq \mathbb Q(\omega_1, \ldots, \omega_{p-1})$ y si $\omega \in \mathbb Q(\omega_1, \ldots, \omega_{p-1})$, entonces $\omega_k = (\omega_1)^k \in \mathbb Q(\omega_1)$. Asi que:
\begin{align*}
  [\mathbb Q(\omega_1):\Q] &= p-1
\end{align*}
Pues $x^{p-1} + x^{p-2} + \dots + x + 1 \in \mathbb Q[x]$  es un polinomio irreducible de grado $p-1$ o también:
\begin{align*}
  \mathbb Q(\omega_1) &= \{
a_0 + a_1 \omega_1 + a_2 \left(\omega_1\right)^2 + \ldots + a_{p-2} \left(\omega_1\right)^{p-2} : a_i \in \mathbb Q\}\\
\end{align*}
tiene dimensión $p-1$ sobre $\mathbb Q$. 
Como ademas $\mathbb Q \subset \mathbb Q(\sqrt[p]{2}) \subseteq L$ y $\mathbb Q \subset \mathbb Q(\omega_1) \subseteq L$, entonces:
\begin{align*}
  p\mid [L:\mathbb Q(\sqrt[p]{2})] \quad \text{y} \quad p-1\mid [L:\mathbb Q(\omega_1)].
\end{align*}
por lo que si $[L:\mathbb Q] = n$, entonces $n \geq p(p-1)$.
Ademas notemos que $\mathbb Q(\sqrt[p]{2}) \not \subseteq \mathbb Q(\omega_1)$, y $\mathbb Q(\omega_1) \not \subseteq \mathbb Q(\sqrt[p]{2})$. Por lo tanto, consideramos $L = \mathbb Q(\sqrt[p]{2}, \omega_1)$, verificamos que en efecto $L$ es el cuerpo de descomposición de $f(x)$, pues:
\begin{itemize}
  \item [$\subseteq$] $\mathbb Q(\sqrt[p]{2}, \omega_1) \subseteq L$, pues $\sqrt[p]{2} \in L$ y $\omega_1 \in L$.
  \item [$\supseteq$] Sea $\alpha \in L$ una raíz de $f(x)$, entonces $\alpha = \sqrt[p]{2} \omega_k$ para algún $k = 0,1,\ldots,p-1$. Por lo tanto $\alpha = \sqrt[p]{2} \omega_k = \sqrt[p]{2} (\omega_1)^k \in \mathbb Q(\sqrt[p]{2}, \omega_1)$, ya que $\mathbb Q(\sqrt[p]{2}, \omega_1)$ contiene a $\sqrt[p]{2}$ y a $\omega_1$.
\end{itemize}
Asi por el teorema de la torre:
\begin{align*}
  [L:\mathbb Q] &= [\mathbb Q(\sqrt[p]{2}, \omega_1):\mathbb Q(\sqrt[p]{2})][\mathbb Q(\sqrt[p]{2}):\mathbb Q]\\
\end{align*}
pero como $\Phi_p(x) = x^{p-1} + x^{p-2} + \ldots + x + 1 \in \mathbb Q(\sqrt[p]{2})[x]$  y $\Phi_p(\omega_1) = 0$, 
entonces $[\mathbb Q(\sqrt[p]{2}, \omega_1):\mathbb Q(\sqrt[p]{2})] \geq p-1$ y $[\mathbb Q(\sqrt[p]{2}):\mathbb Q] = p$, por lo que:
\begin{align*}
  [L:\mathbb Q] = n \leq p(p-1) \quad &\text{y} \quad n\geq p(p-1) \quad \Rightarrow \quad
  n = p(p-1).
\end{align*}
Por lo tanto, el grado de la extensión $L$ sobre $\mathbb Q$ es:
\[
   \boxed{[L:\Q]=p(p-1)}.
\]

% % \[
% % x^5 - 2
% % \;=\;
% % \bigl(x - \sqrt[5]{2}\bigr)\,
% % \Bigl(x^4 + \sqrt[5]{2}\,x^3 + \bigl(\sqrt[5]{2}\bigr)^2 x^2
% %        + \bigl(\sqrt[5]{2}\bigr)^3 x + \bigl(\sqrt[5]{2}\bigr)^4\Bigr)
% % \]

% % \[
% % =\;
% % \prod_{k=0}^{4}\Bigl(x - 2^{1/5}\,e^{2\pi i k/5}\Bigr),
% % \quad
% % e^{2\pi i k/5}\,(k=0,1,2,3,4)\;\text{raíces $5$-ésimas de la unidad.}
% % \]

% % Consecuencia: el \emph{polinomio mínimo} de $\alpha:=\sqrt[p]{2}$ es $f$
% % y
% % \[
% %   [\mathbb Q(\alpha):\Q]=\deg f=p
% %   \quad\text{(def.\ de pol.\ mínimo y
% %         resultado de grado: \textbf{Clase 6.pdf}, p.~2).}
% % \]

% \vspace{0.4em}
% \textbf{2. Raíces de $f$ y descripción del cuerpo de descomposición.}

% Sea $\zeta_{p}=e^{2\pi i/p}$ una raíz $p$-ésima primitiva de la unidad.  
% Las $p$ raíces de $f$ son
% \[
%    \alpha,\;
%    \alpha\zeta_{p},\;
%    \alpha\zeta_{p}^{2},\;
%    \dots,\;
%    \alpha\zeta_{p}^{p-1}.
% \]
% Por tanto
% \[
%    L=\Q\bigl(\alpha,\zeta_{p}\bigr).
% \]

% \vspace{0.4em}
% \textbf{3. Grado de $\mathbb Q(\zeta_{p})/\Q$.}

% El polinomio ciclotómico de orden $p$ es
% \(
%    \Phi_{p}(x)=x^{p-1}+x^{p-2}+\dots+x+1.
% \)
% Usando el truco $x=(y+1)$ y el criterio de Eisenstein
% (ver ejemplo en \textbf{Clase 4.pdf}, p.~3),
% se prueba que $\Phi_{p}(x)$ es irreducible sobre $\Q$
% y $\deg\Phi_{p}=p-1$.  
% Por lo tanto
% \[
%    [\,\mathbb Q(\zeta_{p}):\Q\,]=p-1.
% \]

% \vspace{0.4em}
% \textbf{4. Intersección de los subcampos.}

% Los subcampos intermedios
% \[
%    \mathbb Q(\alpha)\subseteq L, \qquad
%    \mathbb Q(\zeta_{p})\subseteq L
% \]
% tienen grados $p$ y $p-1$, respectivamente.
% Esos números son coprimos, así que
% \[
%     \mathbb Q(\alpha)\cap\mathbb Q(\zeta_{p})=\Q.
% \]
% (Una extensión de grado divisor común $d=\gcd(p,p-1)=1$
%  debe ser la base \emph{per se};
% véase la observación sobre intersecciones coprimas en
% \textbf{Clase 8.pdf}, p.~3).

% \vspace{0.4em}
% \textbf{5. Ley de la torre.}

% Por la \textit{ley de la torre}
% (\textbf{Clase 6.pdf}, p.~2):
% \[
%   [\,L:\Q\,]
%   \;=\;
%   [\,L:\mathbb Q(\zeta_{p})\,]\,[\,\mathbb Q(\zeta_{p}):\Q\,].
% \]
% Pero $L=\mathbb Q(\alpha,\zeta_{p})$ y $\mathbb Q(\alpha)\cap\mathbb Q(\zeta_{p})=\Q$, de modo que
% \[
%    [\,L:\mathbb Q(\zeta_{p})\,]=[\mathbb Q(\alpha):\Q]=p.
% \]
% Por tanto
% \[
%    [L:\Q]=p\,(p-1).
% \]

% \[
%    \boxed{[L:\Q]=p(p-1)}.
% \]
\end{solucion}